% ---------------------------------------------------------------------
%  02-gaps.tex -- where today's methods fall short
% ---------------------------------------------------------------------
\section{Where Today's Methods Fall Short}
\label{sec:gaps}

HPB does not replace the methods below; I still use most of them. It sits
on top of them and fills a specific gap that each one leaves open. This
section goes through them one family at a time, the way they show up in a
product manager's week, and says plainly what each does well and where it
stops. A comparison table and a three-point gap statement close the
section.

\subsection{Backlogs and Kanban boards}
\label{sec:gaps-boards}

A backlog is an ordered list of what should be done next. A Kanban
board~\citep{anderson2010kanban} shows where each item is in the workflow,
limits work in progress (WIP) and makes flow visible. Both are excellent at
answering ``what should the team do now?'' and ``where is it stuck?''.
Reinertsen's work on product development flow~\citep{reinertsen2009flow}
explains why that matters: queues and WIP drive lead time far more than
most people expect, and Little's law~\citep{little1961} ties them
together.

What a board does not record is the \emph{product state}. A card in
``Done'' tells me the work is finished, not which version contains it. A
card in ``In progress'' does not say whether the defect was seen on v1.3 or
v1.4. Two cards for the same product sit next to two cards for another
product, in the same column, with nothing that shows they belong to
different lines. And once the card is archived, the link between the work
and the version it went into is gone. In practice the board answers
\emph{where is the work}, never \emph{where is the product}.

\subsection{Scrum}
\label{sec:gaps-scrum}

Scrum~\citep{schwaber2020scrum} gives a team a rhythm: an ordered product
backlog, time-boxed sprints and a usable increment at the end of each one.
For a software team shipping continuously, the increment and the release
can be almost the same thing.

For a product that combines hardware and firmware they are not. A sprint
increment is not a firmware release that goes through certification, nor a
hardware revision that goes through a pilot lot. Several sprints, sometimes
of several teams, feed one product version; some work lives in the lab and
never enters a sprint at all. Scrum organizes the team's time well, but it
does not say which product baseline an item started from or ended in, and
it was never meant to.

\subsection{Roadmaps and Gantt charts}
\label{sec:gaps-roadmaps}

Roadmaps and Gantt charts answer ``when''. They are the tool management
asks for, and they are good at showing dates, dependencies between
activities and the big picture of a portfolio over quarters.

Their weakness is that they plan activities, not product states. A bar
that says ``low-power mode, March to May'' does not say which version will
carry it, and when the date slips nobody can see which release is now
exposed. They are also forward-looking: once the period is over, a roadmap
is rarely updated to show what actually happened, so it cannot be used for
audit. In my experience the roadmap and the reality drift apart quickly,
and the reconciliation is done in a meeting, not in the artifact.

\subsection{Stage-Gate}
\label{sec:gaps-stagegate}

Stage-Gate~\citep{cooper1990stagegate} brought something very valuable
into product development: explicit decision points. At each gate a
cross-functional review decides, against known criteria, whether to
continue. The phases-and-decisions view of product
development~\citep{ulrich2016pdd} is still how most hardware programs are
run.

The gates, however, apply to \emph{projects}. Once the product is on the
market and evolving through firmware releases, minor hardware revisions
and a steady stream of improvements and fixes, there is no project left to
gate, only a sequence of versions. What I needed was the same discipline
applied to each release: this version does not ship until the work
committed to it is done and verified, or explicitly deferred.

\subsection{PLM and engineering change orders}
\label{sec:gaps-plm}

On the hardware side, PLM systems~\citep{stark2015plm} and engineering
change management~\citep{jarratt2011ec} give strong traceability: change
requests and orders, item revisions, effectivity, approvals. They build on
the configuration management discipline of
baselines, change control and status
accounting~\citep{ieee828,iso10007}, which is very close in spirit to HPB.

The gap is scope and form. Change orders are document-centric and live in
a different system from the tracker where firmware and software defects and
improvements are managed. The PM ends up with hardware lineage in one tool
and the rest of the product's evolution in another, and nothing that shows
both on the same timeline. Nor does a change order give a quick picture of
everything open on a product at once.

\subsection{Git and git-flow}
\label{sec:gaps-git}

Version control is where the right model already exists. In
Git~\citep{chacon2014progit} every line of work has an explicit origin,
parallel lines are first-class, and every merge is a recorded event, so the
content of each release can be derived from the graph. Git-flow~\citep{driessen2010gitflow}
added roles (feature, release, hotfix) and a horizontal diagram that
developers read at a glance; trunk-based development and continuous
delivery~\citep{humble2010cd} pushed integration further. Versioning
conventions such as Semantic Versioning~\citep{semver} even encode the
nature of a change in the version name.

But it all stops at the repository. A git graph does not contain the
customer's defect report, the lab evaluation of a new sensor, the
hardware revision forced by a shortage, or the decision to drop an
improvement because the bill of materials got too expensive. And a product
with firmware, an app and a cloud service has several repositories, none of
which is ``the product''. The model is right; it is applied at the wrong
level.

\subsection{Issue trackers with version fields}
\label{sec:gaps-trackers}

A word on the tool most teams use every day. Issue trackers usually offer
\emph{affects version} and \emph{fix version} fields, and release pages
that list unresolved items. Used with discipline, they record origin and
closure for each item. They still give no structural view across items: no
picture of what is running in parallel on a product, no lineage between
hardware references, and nothing that stops a version from being marked
released while committed items are open. HPB can be built directly on top
of such a tracker (\cref{sec:implementation}).

\subsection{Comparison and the gap}
\label{sec:gaps-comparison}

\Cref{tab:comparison} rates these approaches against eight capabilities
that follow from the questions in \cref{sec:introduction}, listed in its
caption: from recording where work starts and which version closes it, to
hardware lineage, release gates, maintained versions and configurations.
Each approach is rated as it is commonly practiced, not as it could be
extended; most tools can be configured to cover more.

\begin{table}[t]
  \centering
  \small
  \caption{How common approaches cover the eight capabilities a product
    manager needs, as they are usually practiced. (a)~fork baseline
    recorded; (b)~closing version recorded; (c)~parallel work per product
    visible; (d)~hardware lineage preserved; (e)~release gate on committed
    work; (f)~one visual artifact for planning and audit; (g)~several
    maintained versions in parallel; (h)~configurations and variants.}
  \label{tab:comparison}
  \setlength{\tabcolsep}{3pt}
  \begin{tabularx}{\linewidth}{@{}>{\raggedright\arraybackslash}X cccccccc@{}}
    \toprule
    Approach & (a) & (b) & (c) & (d) & (e) & (f) & (g) & (h) \\
    \midrule
    Backlog or list (e.g., Scrum product backlog~\citep{schwaber2020scrum})
      & -- & -- & -- & -- & -- & -- & -- & -- \\
    Kanban board~\citep{anderson2010kanban}
      & -- & -- & $\sim$ & -- & -- & $\sim$ & -- & -- \\
    Roadmap or Gantt chart
      & -- & $\sim$ & $\sim$ & $\sim$ & $\sim$ & $\sim$ & $\sim$ & -- \\
    Tracker with affects/fix version fields
      & \checkmark & \checkmark & $\sim$ & -- & $\sim$ & -- & $\sim$ & $\sim$ \\
    ALM or requirements traceability tool
      & $\sim$ & \checkmark & $\sim$ & -- & $\sim$ & -- & $\sim$ & $\sim$ \\
    Stage-Gate process~\citep{cooper1990stagegate}
      & -- & -- & -- & -- & $\sim$ & -- & -- & -- \\
    PLM engineering change management~\citep{stark2015plm,jarratt2011ec}
      & \checkmark & \checkmark & $\sim$ & \checkmark & \checkmark & -- & $\sim$ & \checkmark \\
    Git with a branching model, for code~\citep{driessen2010gitflow}
      & \checkmark & \checkmark & \checkmark & -- & $\sim$ & $\sim$ & \checkmark & -- \\
    \midrule
    \textbf{HPB (by design)}$^{\dagger}$
      & \checkmark & \checkmark & \checkmark & \checkmark & \checkmark
      & \checkmark & \checkmark & $\sim$ \\
    \bottomrule
  \end{tabularx}

  \smallskip
  \begin{minipage}{\linewidth}
    \footnotesize\emph{Note.} \checkmark: supported as commonly practiced;
    $\sim$: partial, or only through configuration or convention; --: not
    supported. Ratings are my qualitative assessment of typical use, not
    the result of a systematic tool review. Roadmaps get $\sim$ for (d)
    because they often show a successor reference as a new bar, without a
    recorded link. Stage-Gate gets $\sim$ for (e) because its gates apply
    to projects, not to each release of a product in the market. PLM change
    orders gate the release of affected items and record effectivity, but
    they rarely cover firmware and software work and give no overview,
    hence -- for (f). Git models get $\sim$ for (f) because their graph
    covers code only, and $\sim$ for (e) because protected branches and
    checks gate merges, not product releases. $^{\dagger}$HPB is rated by
    construction, not by evidence from use (\cref{sec:limitations});
    configurations are an optional extension and variants are drawn as
    sibling lines, hence $\sim$ for (h).
  \end{minipage}
\end{table}

Read across, the table says two things. First, almost every capability
already exists somewhere: Git has origin, parallelism and closure for
code; trackers record versions item by item; PLM has lineage, gates and
effectivity for hardware; Stage-Gate has decision points; Kanban makes flow
visible. Second, they are almost never available in one place, at the
level of the \emph{product}, across hardware, firmware and software, which
is exactly the level at which a product manager plans releases and is held
accountable for them. HPB does not claim to invent these capabilities. Its
contribution is to put them together in one typed graph, with checkable
rules, a standard picture and an operating routine.

\paragraph{The gap, in three points.}
\begin{itemize}[leftmargin=1.5em,itemsep=2pt,topsep=3pt]
  \item \textbf{No product-level graph.} Origin, parallelism and closure
        are recorded for code, not for the product; work on the product is
        tracked as items, not as branches of a line of versions.
  \item \textbf{Hardware and software live apart.} Hardware lineage sits in
        PLM, firmware and software work in trackers and repositories, and
        nothing shows both on the same timeline.
  \item \textbf{Releases are not gates.} Commitments to a version are
        implicit, so a version can ship while committed work is still open,
        and nobody can later prove what was in it and why.
\end{itemize}

Research on software evolution, traceability and release planning covers
parts of this; its oldest lesson, that a system in use must keep changing
or become less useful~\citep{lehman1980}, is the daily reality of a
product manager. HPB records the link between work and product states
\emph{when each decision is made}, for the whole product, hardware
included.
